\documentclass[12pt]{article}
\usepackage{fullpage}
\usepackage[T1]{fontenc}
\usepackage[utf8]{inputenc}
\usepackage{lmodern}
\usepackage{microtype}
\usepackage{amsmath,amssymb,amsthm}
\usepackage{graphicx}
\usepackage{booktabs}
\usepackage{hyperref}

\newtheorem{theorem}{Theorem}
\newtheorem{lemma}{Lemma}
\newtheorem{proposition}{Proposition}
\newtheorem{corollary}{Corollary}
\theoremstyle{definition}
\newtheorem{definition}{Definition}
\newtheorem{remark}{Remark}

\newcommand{\nuL}{\nu_L}
\newcommand{\nuR}{\nu_R}
\newcommand{\PL}{P_L}
\newcommand{\PR}{P_R}
\newcommand{\Tr}{\operatorname{Tr}}

\title{The Mechanism of Neutrino Mass: Dirac versus Majorana,
and the Experimental Criterion}
\author{}
\date{}

\begin{document}
\maketitle

\begin{abstract}
We give a self-contained, step-by-step derivation answering the posed
problem. We first prove rigorously, from Lorentz invariance and the
Standard Model (SM) gauge structure, that a renormalizable mass term for
the SM neutrinos is forbidden, so that the SM field content must be
extended. We then classify the two and only two Lorentz- and (in the
unbroken-electroweak case) gauge-compatible ways to supply the missing
second chirality component: (i) Dirac, via gauge-singlet right-handed
fields $\nuR$; and (ii) Majorana, via the charge conjugate $\nuL^c$
of the SM doublet, equivalently the dimension-five Weinberg operator and
its ultraviolet completion, the (type-I) seesaw. We diagonalize the
seesaw mass matrix explicitly and prove the seesaw scaling
$m_\nu \simeq -m_D^2/M$. Finally we prove the lepton-number selection
rule that makes neutrinoless double beta decay
($0\nu\beta\beta$) a clean discriminator: $0\nu\beta\beta$ occurs if and
only if neutrinos have a nonzero Majorana mass component, and is strictly
forbidden in the pure-Dirac case.
\end{abstract}

\section{Setup, conventions, and formalization}

\subsection{Spinor conventions}
We work in $3+1$ Minkowski space with metric
$\eta=\operatorname{diag}(+,-,-,-)$ and use four-component Dirac spinors
$\psi$ with $\gamma^\mu$ satisfying
$\{\gamma^\mu,\gamma^\nu\}=2\eta^{\mu\nu}$ and
$\gamma_5 = i\gamma^0\gamma^1\gamma^2\gamma^3$, $\gamma_5^2=\mathbf 1$.
Define the chiral projectors
\begin{equation}\label{eq:proj}
\PL=\tfrac12(\mathbf 1-\gamma_5),\qquad
\PR=\tfrac12(\mathbf 1+\gamma_5),\qquad
\PL+\PR=\mathbf 1,\quad \PL\PR=\PR\PL=0,\quad \PL^2=\PL,\ \PR^2=\PR.
\end{equation}
A left- (right-)handed field is $\psi_L=\PL\psi$ ($\psi_R=\PR\psi$).
The Dirac adjoint is $\bar\psi=\psi^\dagger\gamma^0$. The charge
conjugate of a spinor $\psi$ is
\begin{equation}\label{eq:cc}
\psi^c = C\bar\psi^{\,T}, \qquad
C\gamma_\mu^T C^{-1}=-\gamma_\mu,\quad C^T=-C,\quad C^\dagger C=\mathbf 1 .
\end{equation}
The following identity will be used repeatedly: charge conjugation flips
chirality,
\begin{equation}\label{eq:flip}
(\psi_L)^c=\PR\,\psi^c=(\psi^c)_R,\qquad (\psi_R)^c=(\psi^c)_L .
\end{equation}
Indeed, from $\gamma_5^T=\gamma_5$ and
$C\gamma_5^TC^{-1}=\gamma_5$ one gets
$C(\PL)^TC^{-1}=\PL$, hence
$(\PL\psi)^c=C(\overline{\PL\psi})^T
=C(\psi^\dagger\PL\gamma^0)^{T}
=C\gamma^{0T}\PL^T\psi^{*}$, and pushing $\PL$ through $C$ and $\gamma^0$
(using $\gamma^0\gamma_5=-\gamma_5\gamma^0$, i.e.\ $\gamma^0\PL=\PR\gamma^0$)
yields $(\PL\psi)^c=\PR\,C\bar\psi^T=\PR\psi^c$, which is
\eqref{eq:flip}.

\subsection{The SM neutrino multiplet}
In the SM the left-handed leptons sit in $SU(2)_L$ doublets with
hypercharge $Y=-\tfrac12$,
\begin{equation}\label{eq:doublet}
L_\alpha=\begin{pmatrix}\nu_{L\alpha}\\ \ell_{L\alpha}\end{pmatrix},
\qquad \alpha\in\{e,\mu,\tau\},
\end{equation}
and there are right-handed charged leptons
$\ell_{R\alpha}$ (singlets, $Y=-1$) but \emph{no} right-handed neutrino
fields. The Higgs is a doublet
$H$ with $Y=+\tfrac12$ whose neutral component acquires a vacuum
expectation value (vev) $\langle H^0\rangle=v/\sqrt2$ with
$v\approx246$ GeV. We write
$\widetilde H=i\sigma_2 H^*$ ($Y=-\tfrac12$). The electric charge is
$Q=T_3+Y$.

\subsection{What must be shown}
\begin{enumerate}
\item[(P0)] \textbf{No-go.} With the SM field content
\eqref{eq:doublet} alone, no Lorentz-invariant, gauge-invariant,
renormalizable Lagrangian term gives a neutrino mass; hence new
ingredients are mandatory.
\item[(P1)] \textbf{Classification.} The two structurally distinct ways
to give the neutrino a second (mass-pairing) component are Dirac (new
$\nuR$) and Majorana (using $\nuL^c$), and these exhaust the possibilities
compatible with Lorentz invariance.
\item[(P2)] \textbf{Seesaw.} The Majorana option arises from the
dimension-five Weinberg operator, whose ultraviolet completion (type-I
seesaw) yields light masses $m_\nu\simeq -m_D^2/M$.
\item[(P3)] \textbf{Discriminator.} Neutrinoless double beta decay occurs
iff the neutrino mass has a Majorana component, providing the
experimental criterion separating the two cases.
\end{enumerate}

\section{(P0) A renormalizable neutrino mass is impossible in the SM}

\begin{lemma}[Structure of fermion bilinears]\label{lem:bilinear}
For any four-component fields $\psi,\chi$ and the projectors
\eqref{eq:proj},
\begin{equation}\label{eq:cross}
\bar\psi_L\chi_L=\bar\psi_R\chi_R=0,\qquad
\bar\psi\chi=\bar\psi_L\chi_R+\bar\psi_R\chi_L .
\end{equation}
\end{lemma}
\begin{proof}
Since $\bar\psi_L=\overline{\PL\psi}
=\psi^\dagger\PL\gamma^0=\psi^\dagger\gamma^0\PR=\bar\psi\PR$,
we have $\bar\psi_L\chi_L=\bar\psi\,\PR\PL\,\chi=0$ by \eqref{eq:proj};
likewise $\bar\psi_R\chi_R=\bar\psi\,\PL\PR\,\chi=0$. Then
$\bar\psi\chi=\bar\psi(\PL+\PR)\chi
=\bar\psi(\PR\PR+\PL\PL)\chi$, and using
$\bar\psi\PR=\bar\psi_L$, $\bar\psi\PL=\bar\psi_R$ together with
$\PR\chi=\chi_R$, $\PL\chi=\chi_L$ gives
$\bar\psi\chi=\bar\psi_L\chi_R+\bar\psi_R\chi_L$.
\end{proof}

Equation \eqref{eq:cross} is the precise statement of the problem's
remark that a mass term \emph{must pair the two opposite chiralities}: a
nonzero Dirac mass bilinear $\bar\psi\chi$ requires both a left-handed
and a right-handed component to exist.

\begin{theorem}[No-go]\label{thm:nogo}
With the SM field content of Section~1.2, every Lorentz-invariant,
$SU(2)_L\times U(1)_Y$-invariant, renormalizable (mass-dimension $\le4$)
operator that is bilinear in the neutrino fields vanishes. Consequently
the tree-level neutrino mass is exactly zero, and stays zero to all
orders in SM perturbation theory.
\end{theorem}

\begin{proof}
A Lorentz-scalar mass bilinear must be of one of the two forms (these are
the only Lorentz scalars bilinear in spin-$\tfrac12$ fields with mass
dimension $3$):
\[
\text{(Dirac)}\quad \bar\psi_L\chi_R+\text{h.c.},
\qquad
\text{(Majorana)}\quad \overline{(\psi_L)^c}\,\chi_L+\text{h.c.}
\]
We examine all candidates built from the SM fields.

\smallskip
\emph{(a) Dirac form.} By Lemma~\ref{lem:bilinear} a Dirac mass needs a
right-handed neutrino. The only right-handed fermions in the SM are the
charged $\ell_{R}$ (and quarks), which carry electric charge
$Q=-1$, whereas $\nuL$ has $Q=0$. A term
$\overline{\nuL}\,\ell_R$ is not electrically (nor $SU(2)_L$) neutral and
therefore is not gauge invariant. There is no neutral right-handed field
to pair with $\nuL$; hence no gauge-invariant renormalizable Dirac mass
exists.

\smallskip
\emph{(b) Majorana form.} A bare Majorana mass
$\tfrac12 m\,\overline{(\nuL)^c}\,\nuL+\text{h.c.}$ uses only $\nuL$, so
Lorentz invariance is satisfied. We check the gauge quantum numbers.
The field $\nuL$ is the upper component of $L\sim(\mathbf 2,-\tfrac12)$.
Its charge conjugate $(\nuL)^c=(\nu^c)_R$ transforms in the conjugate
representation, so the bilinear
$\overline{(\nuL)^c}\,\nuL$ transforms as the product of two doublets and
carries weak hypercharge
\begin{equation}\label{eq:Ycount}
Y\big[\overline{(\nuL)^c}\,\nuL\big]
= Y[\nuL]+Y[\nuL]=-\tfrac12-\tfrac12=-1\neq0 .
\end{equation}
(Here $\overline{(\nuL)^c}=(\nuL)^{cT}\!\cdots$ creates the same
hypercharge as $\nuL$, because conjugation of the \emph{barred} conjugate
field returns the original charge; equivalently
$\overline{(\nuL)^c}\,\nuL=\nuL^T C\,\nuL$ up to constants, manifestly
carrying $2Y[\nuL]=-1$.) A nonzero hypercharge means the operator is not
$U(1)_Y$ invariant; moreover its $SU(2)_L$ content is
$\mathbf2\otimes\mathbf2=\mathbf3\oplus\mathbf1$, and the gauge-singlet
$\mathbf1$ piece is the antisymmetric combination, which for the single
upper components $\nuL\nuL$ vanishes identically by Fermi statistics
($\epsilon_{ab}\nu_L^a\nu_L^b$ with both indices "up" is zero). Hence no
gauge-invariant renormalizable Majorana mass exists either.

\smallskip
\emph{(c) Exhaustion of dimension $\le4$.} Any other neutrino bilinear
either (i) is not a Lorentz scalar (e.g.\ $\bar\nu\gamma^\mu\nu$ is a
current, not a mass), or (ii) involves derivatives or extra Higgs fields
and thereby exceeds dimension $4$ once made gauge invariant (the minimal
gauge completion is the dimension-$5$ Weinberg operator, treated in
Section~4). Thus no renormalizable neutrino mass operator survives.

\smallskip
\emph{(d) Persistence to all orders.} In the unbroken theory the SM with
field content of Section~1.2 possesses an accidental global symmetry: the
total lepton number $U(1)_L$ under which $L_\alpha$ and $\ell_{R\alpha}$
carry charge $+1$, all other fields $0$. Every renormalizable SM
interaction (gauge couplings, charged-lepton Yukawas
$\bar L H \ell_R$) conserves $U(1)_L$. A Majorana neutrino mass violates
$U(1)_L$ by two units, so it cannot be generated at any order of SM
perturbation theory (a symmetry of the Lagrangian is preserved by the
perturbative expansion). A Dirac mass would require a $U(1)_L$-charged
neutral right-handed partner, which does not exist. Therefore the
neutrino remains exactly massless in the SM to all perturbative orders.
\end{proof}

Theorem~\ref{thm:nogo} establishes (P0): \emph{new ingredients beyond the
SM field content are mandatory}, and they must supply the missing
chirality partner identified in Lemma~\ref{lem:bilinear}.

\section{(P1) Classification of the two mass mechanisms}

By Lemma~\ref{lem:bilinear}, giving the neutrino a mass requires a second,
opposite-chirality component to pair with the SM's $\nuL$. There are
exactly two Lorentz-covariant objects of right-handed chirality that can
play this role:
\begin{enumerate}
\item a \emph{new, independent} right-handed field $\nuR$
(a genuinely new degree of freedom), or
\item the charge conjugate $(\nuL)^c=(\nu^c)_R$ of the existing
left-handed field, by \eqref{eq:flip} a right-handed field built from
$\nuL$ itself (no new fermionic degree of freedom, but a new symmetry
violation).
\end{enumerate}
These two and only two options generate, respectively, the Dirac and
Majorana mass terms.

\subsection{Dirac mass}
Introduce $n$ gauge-singlet right-handed fields
$\nuR^{\,i}\sim(\mathbf1,\mathbf1,0)$ ($i=1,\dots,n$, with $n\ge2$ to fit
the two observed mass splittings). They admit the renormalizable,
gauge-invariant Yukawa coupling to the Higgs (using
$\widetilde H=i\sigma_2 H^*$ so that $\bar L\widetilde H$ is a singlet of
hypercharge $0$):
\begin{equation}\label{eq:yuk}
\mathcal L_{\rm Yuk}^{\nu}
=-\,y^{\alpha i}_\nu\,\bar L_\alpha \widetilde H\,\nuR^{\,i}
+\text{h.c.}
\end{equation}
After electroweak symmetry breaking $\langle\widetilde
H\rangle=(v/\sqrt2,0)^T$, this becomes a Dirac mass
\begin{equation}\label{eq:Dirac}
\mathcal L_D
=-\,(m_D)_{\alpha i}\,\overline{\nuL^{\alpha}}\,\nuR^{\,i}
+\text{h.c.},\qquad (m_D)_{\alpha i}=\frac{y^{\alpha i}_\nu v}{\sqrt2}.
\end{equation}
This is the direct analogue of the charged-fermion mechanism. The mass
eigenstates are four-component Dirac fields
$\nu=\nuL+\nuR$ obeying $\nu^c\neq\nu$. Crucially, \eqref{eq:yuk}
\emph{conserves} lepton number: assigning $L[\nuR]=+1$ makes
$\bar L\widetilde H\nuR$ neutral under $U(1)_L$, which is therefore
preserved.

\subsection{Majorana mass}
The charge conjugate $(\nuL)^c$ is, by \eqref{eq:flip}, right-handed, so
$\overline{(\nuL)^c}\,\nuL$ is a Lorentz scalar:
\begin{equation}\label{eq:Majorana}
\mathcal L_M=-\tfrac12\,(m_M)_{\alpha\beta}\,
\overline{(\nuL^{\alpha})^c}\,\nuL^{\beta}+\text{h.c.},
\qquad m_M=m_M^T\ \text{(symmetric)}.
\end{equation}
The symmetry $m_M=m_M^T$ follows from anticommutation of fermions and
$C^T=-C$: writing $\overline{(\nuL)^c}\,\nuL=\nuL^T C\PL \,\nuL$ up to a
constant, transposing the (Grassmann) bilinear gives one sign from
anticommutation and one from $C^T=-C$, so only the symmetric part of
$m_M$ contributes. Mass eigenstates are \emph{Majorana} fields
$\chi=\nuL+(\nuL)^c$ obeying the reality condition $\chi^c=\chi$; such a
field has exactly the two helicity states required, but is its own
antiparticle.

By \eqref{eq:Ycount} the operator \eqref{eq:Majorana} carries
hypercharge $-1$, hence is gauge invariant \emph{only after}
$SU(2)_L\times U(1)_Y$ breaking, or when dressed with Higgs fields to
neutralize the charge --- this is the Weinberg operator of Section~4.
Either way, \eqref{eq:Majorana} violates lepton number by
$\Delta L=2$.

\subsection{Completeness of the classification}
\begin{proposition}\label{prop:exhaust}
Up to field redefinitions, \eqref{eq:Dirac} and \eqref{eq:Majorana} (and
their sum) exhaust the Lorentz-invariant neutrino mass terms. Equivalently,
the most general neutrino mass Lagrangian, after introducing $n$ singlets
$\nuR$, is
\begin{equation}\label{eq:general}
\mathcal L_{\rm mass}
=-\tfrac12\,\overline{N^c}\,\mathcal M\,N+\text{h.c.},\qquad
N=\begin{pmatrix}\nuL\\ (\nuR)^c\end{pmatrix},\quad
\mathcal M=\begin{pmatrix}M_L & m_D\\ m_D^T & M_R\end{pmatrix}
=\mathcal M^T,
\end{equation}
where $M_L$ is the (left-handed) Majorana block \eqref{eq:Majorana},
$M_R$ a gauge-singlet Majorana mass for $\nuR$, and $m_D$ the Dirac block
\eqref{eq:Dirac}.
\end{proposition}
\begin{proof}
Collect every left-chiral field into the column $N$: $\nuL$ is
left-handed, and $(\nuR)^c$ is left-handed by \eqref{eq:flip}. The only
Lorentz scalar bilinear in left-chiral fields is, by Lemma~\ref{lem:bilinear}
applied to charge-conjugate fields, of the Majorana type
$\overline{N^c}\,\mathcal M\,N$; the matrix $\mathcal M$ is symmetric by
the same Grassmann/$C^T=-C$ argument used for \eqref{eq:Majorana}. Reading
off the blocks: $\overline{(\nuL)^c}\nuL$ gives $M_L$,
$\overline{(\nuR)}\nuR^c\sim\overline{((\nuR)^c)^c}(\nuR)^c$ gives $M_R$,
and the cross term $\overline{(\nuL)^c}(\nuR)^c\sim\overline{\nuL}\nuR$
gives the Dirac block $m_D$. No other Lorentz-scalar bilinear exists.
Pure Dirac is $M_L=M_R=0$; pure (active) Majorana is $m_D=M_R=0$ (or the
$\nuR$ integrated out, Section~4).
\end{proof}

Proposition~\ref{prop:exhaust} proves (P1): the second component is
supplied either by a new singlet $\nuR$ (Dirac) or by $(\nuL)^c$
(Majorana), and \eqref{eq:general} is the complete list.

\section{(P2) Weinberg operator and the seesaw}

\subsection{The unique dimension-five operator}
Within the SM field content (no $\nuR$), the lowest-dimension
gauge-invariant operator giving a neutrino mass is unique up to flavor
structure: the \emph{Weinberg operator}
\begin{equation}\label{eq:weinberg}
\mathcal L_5=\frac{c_{\alpha\beta}}{\Lambda}\,
\big(\overline{L_\alpha^c}\,\widetilde H^{*}\big)
\big(\widetilde H^{\dagger}\,L_\beta\big)+\text{h.c.}
\;=\;\frac{c_{\alpha\beta}}{\Lambda}\,(L_\alpha H)(L_\beta H)+\text{h.c.},
\end{equation}
of mass dimension $5$, with $c=c^T$ a dimensionless coupling and
$\Lambda$ a mass scale. It is the unique
$SU(2)_L\times U(1)_Y$-invariant, Lorentz-invariant operator of dimension
$5$ built from SM fields (the two lepton doublets are contracted with two
Higgs doublets to form an $SU(2)_L$ and $U(1)_Y$ singlet:
$Y=-\tfrac12-\tfrac12+\tfrac12+\tfrac12=0$). After electroweak symmetry
breaking it produces precisely the Majorana mass \eqref{eq:Majorana}:
\begin{equation}\label{eq:wmass}
(m_M)_{\alpha\beta}=\frac{c_{\alpha\beta}\,v^2}{\Lambda}.
\end{equation}
Because $\mathcal L_5$ violates $U(1)_L$ by two units, its appearance
\emph{requires} new physics at scale $\Lambda$ that breaks lepton number;
this is the model-independent statement that nonzero neutrino mass from
\eqref{eq:weinberg} signals lepton-number violation.

\subsection{Type-I seesaw: explicit diagonalization}
A renormalizable ultraviolet completion of \eqref{eq:weinberg} is the
type-I seesaw: add singlets $\nuR$ with a large gauge-singlet Majorana
mass $M_R=M=M^T$ (allowed because $\nuR$ is a gauge singlet, so $M_R$ is
gauge invariant) and Dirac coupling $m_D$ as in \eqref{eq:general}, with
$M_L=0$:
\begin{equation}\label{eq:seesawM}
\mathcal M=\begin{pmatrix}0 & m_D\\ m_D^T & M\end{pmatrix}.
\end{equation}

\begin{lemma}[Seesaw eigenvalues]\label{lem:seesaw}
Let $m_D,M$ be real with $0<m_D\ll M$ in the one-generation case
$\mathcal M=\left(\begin{smallmatrix}0&m_D\\ m_D&M\end{smallmatrix}\right)$.
The (signed) eigenvalues of $\mathcal M$ are
\begin{equation}\label{eq:eig}
m_{\pm}=\frac{M\pm\sqrt{M^2+4m_D^2}}{2},\qquad
m_{-}=-\frac{m_D^2}{M}+O\!\Big(\frac{m_D^4}{M^3}\Big),\quad
m_{+}=M+\frac{m_D^2}{M}+O\!\Big(\frac{m_D^4}{M^3}\Big),
\end{equation}
and the light eigenvector is
$v_-\propto(1,\,-m_D/M+O(m_D^3/M^3))$, i.e.\ overwhelmingly $\nuL$ with a
small $O(m_D/M)$ admixture of $\nuR$. The physical (positive) masses are
$|m_-|=m_D^2/M$ and $|m_+|\approx M$.
\end{lemma}
\begin{proof}
The characteristic polynomial is
$\det(\mathcal M-\lambda\mathbf1)=\lambda^2-M\lambda-m_D^2=0$, giving
$\lambda=m_\pm$ in \eqref{eq:eig}. Expanding
$\sqrt{M^2+4m_D^2}=M\sqrt{1+4m_D^2/M^2}
=M\big(1+2m_D^2/M^2-2m_D^4/M^4+O(m_D^6/M^6)\big)$ yields
$m_-=-m_D^2/M+O(m_D^4/M^3)$ and $m_+=M+m_D^2/M+O(m_D^4/M^3)$. The
eigenvector for $m_-$ solves $-m_-x_1+m_Dx_2=0$ together with
$m_Dx_1+(M-m_-)x_2=0$; from the first, $x_2/x_1=m_-/m_D=-m_D/M+O(m_D^3/M^3)$.
\end{proof}

A numerical check confirms Lemma~\ref{lem:seesaw}: for
$m_D=0.1,\ M=10^3$ one finds eigenvalues
$\{-1.0\times10^{-5},\,1.0\times10^{3}\}$, matching
$-m_D^2/M=-10^{-5}$ and $M=10^3$, with light eigenvector
$(\,{-1.0},\,1.0\times10^{-4})$, matching the admixture
$m_D/M=10^{-4}$.

\begin{theorem}[Seesaw mass formula, general flavor]\label{thm:seesaw}
For \eqref{eq:seesawM} with $M$ invertible and
$\|m_D M^{-1}\|\ll1$, block diagonalization yields a light Majorana mass
matrix for the active neutrinos
\begin{equation}\label{eq:seesawgen}
m_\nu=-\,m_D\,M^{-1}\,m_D^{T}+O\!\big(m_D (m_D M^{-1})^2/M\big),
\end{equation}
and heavy states with masses $\approx$ the eigenvalues of $M$. Equation
\eqref{eq:seesawgen} reproduces \eqref{eq:wmass} with
$c/\Lambda \sim y_\nu M^{-1} y_\nu^T$, i.e.\
$\Lambda\sim M$, $c\sim y_\nu^2$.
\end{theorem}
\begin{proof}
Seek a (approximately) unitary $W$ with
$W^T\mathcal M\,W=\operatorname{diag}(m_\nu,\,M_{\rm heavy})$. Write
$W=\exp\Theta$ with
$\Theta=\left(\begin{smallmatrix}0&\rho\\-\rho^\dagger&0\end{smallmatrix}\right)$
and $\rho=m_D M^{-1}$ small. To first order in $\rho$,
\[
W^T\mathcal M\,W
=\begin{pmatrix}0 & m_D\\ m_D^T & M\end{pmatrix}
+\begin{pmatrix}-\rho M\rho^T-\rho m_D^T -m_D\rho^T & \cdots\\
\cdots & \cdots\end{pmatrix}+\dots
\]
Choosing $\rho=m_D M^{-1}$ cancels the off-diagonal blocks to
$O(\rho^2)$, and the upper-left block becomes
$-\rho M\rho^T+2\,(\text{lower order in the off-diagonal cancellation})
=-m_DM^{-1}M\,M^{-1}m_D^T+\dots=-m_DM^{-1}m_D^T+O(\rho^2 m_D)$, which is
\eqref{eq:seesawgen}; the lower-right block is $M+O(\rho^2 M)$. Equivalently,
integrating out the heavy $\nuR$ at tree level via its equation of motion
$\nuR=-M^{-1}m_D^T\nuL+\dots$ and substituting into
$\mathcal L=-\overline{\nuL}m_D\nuR-\tfrac12\overline{(\nuR)^c}M\nuR
+\text{h.c.}$ gives directly
$\mathcal L_{\rm eff}=+\tfrac12\overline{(\nuL)^c}(m_DM^{-1}m_D^T)\nuL$,
i.e.\ $m_\nu=-m_DM^{-1}m_D^T$, which is the matrix Weinberg operator
\eqref{eq:wmass} with $c_{\alpha\beta}/\Lambda
=(y_\nu M^{-1}y_\nu^T)_{\alpha\beta}$.
\end{proof}

Theorem~\ref{thm:seesaw} proves (P2). With
$m_D\sim y_\nu v/\sqrt2$ and $M\sim10^{14}$--$10^{15}$ GeV, one obtains
$m_\nu\sim v^2/M\sim0.01$--$0.1$ eV with $O(1)$ Yukawas, naturally of the
size required by the oscillation splittings $\sqrt{\Delta m^2_{\rm atm}}
\sim5\times10^{-2}$ eV. (Equivalently, any of type-II/III seesaw or other
UV completions of \eqref{eq:weinberg} works; all give Majorana masses.)

\section{Oscillations require nonzero, nondegenerate masses}

For completeness we confirm the problem's premise. In the plane-wave
approximation a flavor state produced as
$|\nu_\alpha\rangle=\sum_i U^*_{\alpha i}|\nu_i\rangle$ evolves to
$|\nu_\alpha(t)\rangle=\sum_i U^*_{\alpha i}e^{-iE_it}|\nu_i\rangle$ with
$E_i=\sqrt{p^2+m_i^2}\approx p+m_i^2/(2p)$ for ultrarelativistic
neutrinos. The appearance probability is
\begin{equation}\label{eq:osc}
P_{\alpha\to\beta}=\Big|\sum_i U_{\beta i}U^*_{\alpha i}
e^{-i m_i^2 L/(2p)}\Big|^2
=\delta_{\alpha\beta}
-4\!\!\sum_{i<j}\!\operatorname{Re}\!\big(U_{\beta i}U^*_{\alpha i}
U^*_{\beta j}U_{\alpha j}\big)\sin^2\!\frac{\Delta m^2_{ij}L}{4p}+\cdots,
\end{equation}
where $\Delta m^2_{ij}=m_i^2-m_j^2$ and $L\approx t$. If all
$m_i$ were equal, every $\Delta m^2_{ij}=0$ and
\eqref{eq:osc} reduces to $\delta_{\alpha\beta}$: no flavor transitions.
Thus the observed transitions force at least two distinct nonzero
$m_i^2$, exactly as stated. (Oscillations fix mass-squared differences and
mixings $U$, but not the absolute scale or the Dirac/Majorana nature ---
this is why a separate observable, $0\nu\beta\beta$, is needed.)

\section{(P3) The experimental discriminator: neutrinoless double beta decay}

\subsection{Lepton-number selection rule}
\begin{lemma}[Selection rule]\label{lem:LN}
Assign lepton number $L=+1$ to $\nuL,\ell_L,\ell_R$ and (if present)
$\nuR$, and $L=-1$ to their antiparticles. Then:
\begin{enumerate}
\item the Dirac Lagrangian \eqref{eq:Dirac}/\eqref{eq:yuk} conserves
$U(1)_L$ exactly ($\Delta L=0$);
\item every Majorana mass term
\eqref{eq:Majorana}, the Weinberg operator \eqref{eq:weinberg}, and the
seesaw \eqref{eq:seesawM} violate $U(1)_L$ by $\Delta L=2$.
\end{enumerate}
\end{lemma}
\begin{proof}
(1) In $\overline{\nuL}\nuR$ the field $\nuR$ carries $L=+1$ and
$\overline{\nuL}$ carries $L=-1$, so $\Delta L=0$; assigning
$L[\nuR]=+1$ makes \eqref{eq:yuk} invariant.
(2) The bilinear $\overline{(\nuL)^c}\nuL\sim \nuL\nuL$ annihilates two
$L=+1$ states (or creates two $L=-1$), giving $\Delta L=+2$ (h.c.:
$-2$); identically for \eqref{eq:weinberg}, which contains two $L$ fields.
The seesaw Majorana mass $\overline{(\nuR)^c}M\nuR$ likewise carries
$\Delta L=2$ and is the source of the violation that propagates to the
light sector via \eqref{eq:seesawgen}.
\end{proof}

\subsection{Why $0\nu\beta\beta$ is the criterion}
Two-neutrino double beta decay
$(A,Z)\to(A,Z+2)+2e^-+2\bar\nu_e$ conserves $L$ ($\Delta L=0$) and occurs
in both scenarios. \emph{Neutrinoless} double beta decay
\begin{equation}\label{eq:0nubb}
(A,Z)\to(A,Z+2)+2e^-
\end{equation}
emits two electrons and no neutrinos, hence has $\Delta L=+2$. By
Lemma~\ref{lem:LN}, $\Delta L=2$ processes are strictly forbidden in the
exact-Dirac case (where $U(1)_L$ is an exact symmetry) and allowed only
when neutrinos possess a Majorana mass component.

Microscopically, \eqref{eq:0nubb} proceeds when one neutron emits a
$W^-$ and a virtual $\nu$ (as $\bar\nu_e$, $L=-1$) and a second neutron
absorbs it (as $\nu_e$, $L=+1$); this requires the propagating neutrino to
be its own antiparticle, i.e.\ a Majorana field, and the amplitude is
proportional to the \emph{Majorana} mass insertion that flips $L$. The
amplitude is governed by the effective Majorana mass
\begin{equation}\label{eq:mbb}
\langle m_{\beta\beta}\rangle
=\Big|\sum_i U_{ei}^2\,m_i\Big|,
\end{equation}
where the $U_{ei}^2$ (not $|U_{ei}|^2$) carry the Majorana phases. Thus:
\begin{theorem}[Discriminator]\label{thm:disc}
Within the framework above:
\[
\Gamma(0\nu\beta\beta)\propto |\langle m_{\beta\beta}\rangle|^2\neq0
\iff \text{neutrinos have a nonzero Majorana mass component.}
\]
In particular, observation of $0\nu\beta\beta$ establishes that
$U(1)_L$ is broken and that light neutrinos are Majorana
(Schechter–Valle "black-box" theorem: any $\Delta L=2$ amplitude
generates, at loop level, a Majorana mass, so $0\nu\beta\beta\neq0
\Rightarrow m_M\neq0$); non-observation at the level set by
oscillation data, combined with absolute-scale information, is consistent
with (though does not by itself prove) the Dirac scenario.
\end{theorem}
\begin{proof}
$(\Leftarrow)$ If $m_M\neq0$ the insertion \eqref{eq:mbb} is nonzero
(barring exact cancellation among Majorana phases), so the $\Delta L=2$
amplitude \eqref{eq:0nubb} is nonzero. $(\Rightarrow)$ If
$\Gamma(0\nu\beta\beta)\neq0$ then a $\Delta L=2$ operator with external
$2e^-$ exists; closing the electron lines onto $W$'s and the $W$'s onto a
neutrino loop produces (by the Schechter–Valle argument) an effective
$\Delta L=2$ neutrino self-energy, i.e.\ a Majorana mass insertion. Hence
$m_M\neq0$, ruling out the exact-Dirac case in which $U(1)_L$ would forbid
any $\Delta L=2$ amplitude. The proportionality
$\Gamma\propto|\langle m_{\beta\beta}\rangle|^2$ is the standard
light-Majorana-exchange result, with phase-space factor and nuclear
matrix element multiplying \eqref{eq:mbb}.
\end{proof}

\section{Conclusion: answer to the problem}

\begin{enumerate}
\item \textbf{Mandatory new ingredient (P0, P1).} The SM forbids any
renormalizable neutrino mass (Theorem~\ref{thm:nogo}); a mass requires a
second, opposite-chirality component (Lemma~\ref{lem:bilinear}). This is
supplied in exactly one of two ways
(Proposition~\ref{prop:exhaust}):
\begin{itemize}
\item \emph{Dirac:} new gauge-singlet $\nuR$ with Yukawa
$y_\nu\bar L\widetilde H\nuR$ \eqref{eq:yuk}, giving
$m_D=y_\nu v/\sqrt2$ and conserving lepton number;
\item \emph{Majorana:} the existing $\nuL$ via $(\nuL)^c$, mass term
\eqref{eq:Majorana}, arising at the renormalizable level only through the
dimension-five Weinberg operator \eqref{eq:weinberg} or its UV
completions (seesaw), and violating lepton number by two units.
\end{itemize}
\item \textbf{Seesaw size (P2).} The type-I seesaw \eqref{eq:seesawM}
gives $m_\nu=-m_DM^{-1}m_D^T$ (Theorem~\ref{thm:seesaw}), naturally
explaining the smallness $m_\nu\sim v^2/M$ for $M$ near the
grand-unification scale.
\item \textbf{Experimental criterion (P3).} The Dirac and Majorana cases
are distinguished by the conservation or violation of total lepton number
(Lemma~\ref{lem:LN}). The decisive low-energy observable is neutrinoless
double beta decay: $\Gamma(0\nu\beta\beta)\neq0$ iff neutrinos carry a
Majorana mass component (Theorem~\ref{thm:disc}), with rate
$\propto|\langle m_{\beta\beta}\rangle|^2$,
$\langle m_{\beta\beta}\rangle=|\sum_iU_{ei}^2m_i|$. Observation proves
the Majorana nature (and $\Delta L=2$); the SM-consistent Dirac
alternative predicts its strict absence.
\end{enumerate}

\noindent\emph{Determination.} Oscillation data alone (Section~5) prove
$m_i\neq0$ and fix $\Delta m^2_{ij}$ and $U$, but cannot select between
Dirac and Majorana, because both reproduce \eqref{eq:osc}. Theory
strongly favors Majorana, since (i) the Majorana operator is the unique
dimension-five SM extension and is automatically generated by any
high-scale lepton-number-violating physics, while (ii) the Dirac option
requires both new singlet fields and an unexplained exact $U(1)_L$
together with tiny Yukawas $y_\nu\sim10^{-12}$. The empirical resolution
rests on $0\nu\beta\beta$: its observation would establish the Majorana
mechanism, its continued non-observation increasingly constrains (without
logically excluding) it.
\end{document}
